
\textbf{7.11.} On peut, pour un schéma \(T\) sur un corps \(k\) algébriquement clos, définir le sous-groupe de \(K^\bullet(T)\) des éléments \emph{algébriquement équivalents à zéro} comme le sous-groupe des éléments de la forme
\begin{equation}\tag{7.11.1}
  E(x_1)-E(x_0),
\end{equation}
où \(x_0\) et \(x_1\) sont deux points \(k\)-rationnels d'une courbe algébrique \(C\) propre, régulière et connexe sur \(k\), \(E\) un élément de \(K^\bullet(T\times_k C)\), et où pour tout point \(x\) rationnel sur \(k\) de \(C\), on pose
\[
  E(x)=i_x^*(E)\in K^\bullet(T),\qquad\text{où }
  i_x:T\longrightarrow T\times_k C
\]
est le morphisme \(t\mapsto(t,x)\) défini par \(x\). Le fait que les éléments de la forme (7.11.1) forment un sous-groupe de \(K^\bullet(T)\) se voit par un raisonnement standard, utilisant un produit \(C\times_k C'\) et le couple de points \((x_1,x'_1)\), \((x_0,x'_0)\) de celui-ci, et le fait que par deux points d'un schéma algébrique quasi-projectif irréductible sur \(k\) (ici \(C\times_k C'\)) on peut toujours faire passer une courbe algébrique irréductible. (Le lecteur qui ne voudra pas admettre ce fait pourra définir l'équivalence algébrique par le groupe engendré par les éléments de la forme (7.11.1)). Les éléments algébriquement équivalents à zéro forment un idéal de \(K^\bullet(T)\), comme on constate aussitôt, d'ailleurs stable par les opérations \(\lambda^i\) lorsque celles-ci sont définies, i.e. lorsque l'on travaille avec le \(K^\bullet\) «~naïf~»~; cet idéal est d'ailleurs contenu dans l'idéal d'augmentation de \(K^\bullet(T)\), comme on vérifie trivialement. On peut donc définir un anneau quotient
\begin{equation}\tag{7.11.2}
  K^\bullet_{\mathrm{alg}}(T),
\end{equation}
qui est même un \(\lambda\)-anneau lorsqu'on travaille avec les \(K^\bullet\) naïfs.

On peut définir de la même façon un sous-groupe des éléments algébriquement équivalents à 0 dans \(K_\bullet(T)\), en prenant ci-dessus \(E\) dans \(K_\bullet(T\times_k C)\), et
